\subsection{Время: почему не sleep}

Чтобы цвет менялся, нужно время. Первое, что приходит в голову, --- остановить программу на секунду функцией \verb|sleep|:

\begin{codefile}{lua/05_leds/leds_sleep.lua}
\end{codefile}

Запустите: красный, через секунду зелёный, ещё через секунду линейка гаснет. Работает. Но на борту Пионера, пока \verb|sleep| «спит», программа \hyperlink{term:blocking}{не делает ничего другого}: не принимает события автопилота, не обслуживает другие таймеры, не читает датчики. На земле это терпимо, в полёте --- нет: дрон не узнает, что долетел до точки, и не отреагирует на разряд батареи.

Симулятор здесь мягче настоящего дрона: его \verb|sleep| приостанавливает только основной текст программы, а события и таймеры продолжают приходить. Поэтому программа со \verb|sleep| в симуляторе может вести себя «правильно» и сломаться на борту. Не полагайтесь на эту мягкость.

\begin{rulebox}[Без sleep]
В программах для дрона \verb|sleep| не используется. Пауза --- это таймер: «через $t$ секунд сделай вот это», а программа тем временем свободна.
\end{rulebox}

\begin{taskbox}[Пауза делает цвет видимым]
Уберите из \texttt{leds\_sleep.lua} обе строки \verb|sleep(1)| и запустите. Затем верните их и добавьте между зелёным и выключением жёлтый, тоже на секунду.
\expected{без пауз линейка сразу выключена: команда цвета не задаёт время показа, и глаз видит только последнее состояние. С паузами --- красный, зелёный, жёлтый по секунде, затем темнота; всего около трёх секунд.}
\mission{Lua}{28}
\end{taskbox}

\subsection{Таймеры}

\begin{codebox}[API таймеров]
\begin{itemize}
  \item \verb|Timer.callLater(delay, fn)| --- вызвать \verb|fn| один раз через \verb|delay| секунд.
  \item \verb|t = Timer.new(period, fn)| --- создать периодический таймер: \verb|fn| вызывается каждые \verb|period| секунд, пока таймер запущен.
  \item \verb|t:start()| и \verb|t:stop()| --- запустить и остановить периодический таймер. Сам по себе \verb|Timer.new| таймер не запускает: без \verb|t:start()| функция не вызовется ни разу.
\end{itemize}
Время всегда в секундах: \verb|0.5| --- полсекунды.
\end{codebox}

\subsubsection*{Отложенное действие}

\begin{codefile}{lua/05_leds/leds_timers_delayed.lua}
\end{codefile}

Программа заканчивается сразу, но через 3 секунды автопилот вызывает переданную функцию, и линейка загорается синим. Обратите внимание: в \verb|callLater| передаётся \emph{функция}, а не результат её вызова.

\subsubsection*{Периодическое действие: мигание}

\begin{codefile}{lua/05_leds/leds_timers_blink.lua}
\end{codefile}

\paragraph{Что происходит.} Переменная \verb|on| хранит \hyperlink{term:state}{состояние} мигалки. Каждые 0.5~с функция \verb|updateBlink| сначала выставляет цвет по текущему состоянию (\verb|on and 1 or 0| --- единица, если \verb|on| истинно), а затем меняет состояние на противоположное (\verb|on = not on|) --- на следующем тике цвет будет другим. Первый тик линейку гасит, второй зажигает красным, и так по кругу. Таймер \verb|blinkTimer| объявлен глобально --- так он живёт до конца программы, и его можно остановить из любой функции.

\begin{hintbox}[Один тик --- один шаг]
Функция таймера должна делать один короткий шаг анимации и возвращаться. Никаких циклов ожидания внутри: следующий шаг таймер вызовет сам.
\end{hintbox}

\begin{taskbox}[Отложить, повторить, пробежать]
\begin{enumerate}
  \item Через 1~с после запуска включите светодиод 0 зелёным и напечатайте \verb|"зелёный"| --- обе строки внутри функции, переданной в \verb|callLater|.
  \item Через 2~с после запуска включите зелёный, через 4~с --- выключите. Оба \verb|callLater| создайте сразу при старте.
  \item Мигание: светлая часть 0.5~с, тёмная 0.5~с. Каков период полного повторения? Затем измените мигалку из примера на синюю с периодом таймера 0.2~с.
  \item «Бегущий огонёк»: таймер с периодом 0.1~с зажигает ровно один зелёный светодиод, индекс которого растёт на 1 и после 28 возвращается к 0.
\end{enumerate}
\expected{1) через секунду зелёный и строка в журнале одновременно; 2) зелёный горит с 2-й по 4-ю секунду: второй таймер считает время от старта, а не от первого; 3) период --- 1~с (свет плюс темнота); синий меняется пять раз в секунду; 4) огонёк обегает всю линейку за 2.9~с, одновременно горит только один светодиод.}
\mission{Lua}{29--31}
\end{taskbox}
